\section{Foundations: The Static Market Environment}
\label{sec:foundations}

This section and Section~\ref{sec:dynamics} establish the
theoretical and quantitative anchor of the thesis: a calibrated
oligopoly model and the classical equilibrium concepts against
which the original contribution is measured. That contribution
itself, the multi-agent learning experiment and its
identification of a path-dependent collusive attractor, is
developed in Section~\ref{sec:marl}. A reader already familiar
with static and repeated oligopoly theory may read the present
two sections as calibration and proceed directly to the learning
experiments; a reader new to the field will find each concept
defined before it is used.

\subsection{Model overview and calibration}
\label{subsec:model}

We model the global crude oil market as a strategic oligopoly. The word
``oligopoly'' denotes an industry in which a small number of producers
each control enough output to influence the market price
\citep{bresnahan1989}; the implication is that production decisions are
inherently strategic, because every producer must anticipate how its
choice affects the price that its rivals also receive. The three blocs
we model are the OPEC core (treated as a single Gulf-led player), the
Russian Federation, and the United States shale sector. Rather than
relying on a single game-theoretic frame, we construct a progressive
computational pipeline in which each layer relaxes one assumption of the
preceding one. The simplest model, static Cournot competition, provides
the non-cooperative benchmark. Every subsequent extension adds one
dimension of realism: timing (Stackelberg), commitment (capacity),
welfare (deadweight loss and carbon taxation), strategic variable
(Bertrand), repeated interaction (Folk Theorem), mediated recommendations
(correlated equilibrium), bargaining (Shapley values), uncertainty
\citep{greenporter1984}, and population-level selection (evolutionary
dynamics). This architecture is consistent with the escalation logic
recommended by \citet[p.~??]{tirole1988} and parallels the framework adopted by
\citet{alcelay2017}, to whose three-player Cournot setting our
calibration is closest.

The demand side of the model is given by a linear inverse demand,
introduced for analytical tractability and because it is the workhorse
specification of the applied oligopoly literature since
\citet[p.~??]{cournot1838}. The function takes the form
\begin{equation}
P(Q) \;=\; a - b\,Q,
\label{eq:1}
\end{equation}
where $P$ denotes the market-clearing price in dollars per barrel,
$Q$ is total output in millions of barrels per day (mbd), $a$ is the
demand intercept (the choke price at which demand falls to zero), and
$b$ is the demand slope (the dollars-per-barrel price response to a
one-unit change in total supply). Each producer
$i \in \{\text{US, OPEC, RUS}\}$ chooses an output $q_i$ to maximise its
period profit
\begin{equation}
\pi_i \;=\; (P - c_i)\, q_i \;=\; (a - b Q - c_i)\, q_i,
\label{eq:2}
\end{equation}
where $c_i$ is its constant marginal cost. The first term in (\ref{eq:2})
is the per-barrel margin and the second is the volume; the product is
the producer's profit-relevant cash flow per period. We calibrate
$a = 140$, $b = 1.0$, $c_{\text{US}} = 45$, $c_{\text{OPEC}} = 20$,
$c_{\text{RUS}} = 35$, and discount factor $\delta = 0.95$. The
calibration is reported in \appref{app:calibration} at the back of the document.

The demand intercept $a = 140$ corresponds to the upper end of the
long-run price scenarios published by the International Energy Agency
\citep{iea2019}; reading it as a choke price is consistent with
high-stress scenarios used by central banks for inflation modelling. The
slope $b = 1.0$ is the standard normalisation in the applied oligopoly
literature \citep[p.~??]{tirole1988}; at the baseline triopoly equilibrium it
implies an own-price demand elasticity of approximately $-0.75$, broadly
consistent with the medium-run elasticities of $-0.4$ to $-1.0$ reported
by the U.S.\ Energy Information Administration in its Short-Term Energy
Outlook. The quarterly discount factor $\delta = 0.95$ implies an
annualized time preference of roughly 19 percent; we read this as a
reduced-form proxy for the political, fiscal, and sanctions-related
impatience that depresses the effective planning horizon of producing
states below the risk-free rate.

We now explain why each marginal cost was set at its calibrated level.
For OPEC we set $c_{\text{OPEC}} = 20$. The Gulf core (Saudi Arabia,
Kuwait, the United Arab Emirates, Iraq) operates onshore conventional
fields with some of the lowest lifting costs in the world: Saudi
Aramco's reported lifting cost has been in the \$3 to \$5 per barrel
range for over a decade, with broader Gulf averages in the \$5 to \$10
range \citep{bp2022,eia2024}. Layering on royalties, infrastructure
depreciation, and the implicit fiscal contribution that funds the state
operating model brings the effective short-run marginal cost to roughly
\$15 to \$20. We take the upper bound of that band to reflect the
average OPEC core producer rather than only Saudi Aramco, and to avoid
overstating the cost advantage in the cartel analysis.

For Russia we set $c_{\text{RUS}} = 35$. Western Siberian conventional
fields have lifting costs of approximately \$15 to \$25 per barrel
\citep[p.~??]{henderson2015,henderson2011}\citep{imf2021}. To this we add the Urals export transport
differential (pipeline tariffs to Baltic and Black Sea ports, plus the
structural discount of Urals to Brent), which contributes a further
\$8 to \$12, and the incremental post-2022 sanctions costs (shadow
fleet, intermediated trade, insurance) of roughly \$5
\citep{simolasolanko2017,eia2025}. The sum sits around \$35, consistent
with the IMF Russia Article~IV estimates \citep{imf2021} and with the
average Urals realization observed in 2023 to 2025. Our model holds this number
constant; in reality, the sanctions overhead has varied with the G7
price cap regime, but the static specification cannot accommodate that
dynamic.

For the United States we set $c_{\text{US}} = 45$. US shale producers
have well-level break-evens that vary widely by play and well quality,
from below \$40 in the core Permian to above \$70 in marginal Bakken
locations \citep{dallasfed2026,smith2016}. We choose \$45 as the average
well-level marginal cost of new wells in major plays, which is the
relevant number for short-run best-response decisions. Firm-level
full-cycle break-evens, which include corporate overhead,
decline-replacement capex, and required returns to equity, were reported
at \$61 to \$62 per barrel in the Q1 2026 Dallas Fed survey
\citep{dallasfed2026}, but those lie outside our short-run analysis. The
post-2020 capital-discipline regime documented by \citet{kraneagerton2015}
and reaffirmed by the same survey implies that US shale's effective supply curve has
flattened, which our static model captures by treating $c_{\text{US}}$
as constant.

Three limitations of the calibration deserve acknowledgment. First,
linear demand abstracts from the documented convexity of crude oil
prices when spare capacity is tight \citep{till2016}. Second, each
``player'' is in reality a coalition: the OPEC core averages over twelve
heterogeneous members, US shale is a fragmented private-sector industry,
and Russia's effective cost varies by destination market
\citep[p.~??]{henderson2015}\citep{eia2025}. Third, marginal costs are held constant,
while in reality they vary with cumulative extraction, technological
learning, and logistics. These simplifications follow standard practice
in applied oligopoly modelling
\citep{salant1976,huppmannholz2012,alcelay2017} and are necessary to
maintain comparability across the frameworks we implement. The full
step-by-step ``escalation table'' is provided in \appref{app:reading-guide}; each row
identifies one framework, the question it addresses, and the assumption
it relaxes relative to its predecessor.

\subsection{Static Cournot equilibrium}
\label{subsec:cournot}

We begin with the simultaneous quantity game, named after
\citet[p.~??]{cournot1838}, the first author to formalize oligopolistic
interaction. In a Cournot game, each producer chooses an output level
without observing the choices of its rivals, and the market-clearing
price is determined by the sum of all outputs through the inverse-demand
function. The Nash equilibrium of a Cournot game is the mutually
best-responding profile: each player's quantity is optimal given the
quantities chosen by everyone else, so no player can unilaterally
improve its profit by deviating. The Cournot frame is the natural
benchmark for oil markets because oil producers actually compete on
quantities, not on posted prices: OPEC announces production targets, not
barrel prices, and the spot price is set in real time by exchange-based
trading on the resulting aggregate supply. We solve the model for both
the duopoly (OPEC and US) and the triopoly (OPEC, US, Russia), following
the three-player structure of \citet{alcelay2017}.

In the duopoly we obtain total output $Q = 71.67$ mbd, equilibrium price
$P = \$68.33/\text{bbl}$, and profits of 2,336.11 for OPEC and 544.44
for US shale.\footnote{Profit, welfare, and surplus figures throughout
are deterministic outputs of the calibrated simulator, reported to two
decimals for reproducibility. The precision reflects the model's
internal arithmetic on stylised parameters, not a claim of empirical
accuracy at that resolution.} Adding Russia as a third strategic player shifts the
equilibrium to $Q = 80.00$ mbd, $P = \$60.00/\text{bbl}$, with profits
of 1,600.00 for OPEC, 225.00 for US, and 625.00 for Russia, consumer
surplus of 3,200.00, and total welfare of 5,650.00 (see Figures~1,
2 and~3). Russia's entry adds 8.33 mbd of output, lowers the price by
\$8.33, erodes OPEC profit by 736.11 units (a 31.5 percent decline), and
raises welfare by 201.39 units. The mechanism is exactly the one
identified by \citet{salant1976} and \citet{novshek1985}: a new entrant
with intermediate cost simultaneously expands aggregate supply and
compresses the incumbent's margin. \citet{schmalensee1987} shows further
that cost asymmetry shifts the unconstrained collusive optimum toward
the low-cost producer, which is why OPEC retains the largest share even
after Russia's entry. This quantitative finding gives a clean economic
reading of the 2016 OPEC+ formation: Russia's standalone presence as a
Cournot rival destroyed roughly 736 profit units for OPEC, which is
precisely the surplus that the OPEC+ coordination mechanism subsequently
attempted to recover.

Comparative statics on the demand intercept confirm the linearity of the
triopoly solution. Sweeping $a$ from 112 to 168, equilibrium price moves
from \$53 to \$67 (slope $dP/da = 1/4$) and total output from 59 to 101
mbd (slope $dQ/da = 3/4$), exactly reproducing the closed-form
predictions of the asymmetric Cournot model with $N = 3$, namely
$dP/da = 1/(N+1)$ and $dQ/da = N/(N+1)$ (see Figure~4). This perfect
linearity is convenient for diagnostics, because any deviation in the
simulator output would signal a coding or calibration error; it is also
the chief reason we are confident in the comparative statics that
follow. The linear response is, however, a property of the demand
specification and should not be over-interpreted. The empirical
price-quantity relationship in physical crude markets is markedly
convex once OPEC spare capacity falls below roughly 2 mbd:
\citet{till2016} documents exactly this regime, in which a small
additional supply deficit produces a disproportionately large price
spike. Our linear model captures the first-order direction of these
effects but not the curvature, and we therefore read the slopes above
as lower bounds on the price sensitivity that would obtain in a
tight-capacity scenario such as the February 2026 Hormuz closure.

\subsection{Market power and first-mover advantage}
\label{subsec:marketpower}

To quantify the degree of market concentration we use two standard
metrics of industrial organisation, both of which we introduce before
reporting our results. The Lerner index $L_i = (P - c_i)/P$, due to
\citet{lerner1934}, measures the gap between price and the marginal cost
of producer~$i$, expressed as a fraction of the price. It equals zero
under perfect competition (price equals marginal cost) and approaches
one as the producer's pricing power rises \citep{bresnahan1989}. The
Herfindahl--Hirschman Index $\mathrm{HHI} = \sum_i s_i^2 \times 10{,}000$,
where $s_i$ is the output share of producer~$i$, measures aggregate
market concentration. It ranges from close to zero in an atomistic
market to 10,000 in a pure monopoly. By the threshold used in the U.S.\
Horizontal Merger Guidelines, an HHI above 2,500 places a market in the
``highly concentrated'' category. Both indices were popularised in the
empirical industrial-organisation literature surveyed in
\citet{bresnahan1989}.

Our simulations yield an HHI above 2,500 in every market structure we
test, which formally establishes that the oil market is, by U.S.\
Department of Justice criteria, highly concentrated; this is not a
trivial finding because the result holds even in the triopoly, when
total industry output is at its largest. The OPEC Lerner index ranges
from 0.57 (when OPEC is the Stackelberg leader, defined below) to 0.71
(in the duopoly), and is the highest among the three players in every
configuration (see Figure~5). The result is consistent with the
calibrated OPEC market-power trajectories of \citet{huppmannholz2012}
and \citet{huppmann2013}, who estimate a conjectural-variation parameter
$\theta$ for OPEC in the range 0.5 to 0.7 throughout 2003 to 2011. In
the standard \citet{bresnahan1989} identity, the industry-weighted
Lerner index satisfies $L = \theta/\varepsilon$; at our equilibrium
elasticity $|\varepsilon| \approx 0.75$, that $\theta$-range translates
to $L \in [0.67, 0.93]$. Our triopoly-Nash OPEC Lerner of
$(60-20)/60 = 0.67$ sits exactly at the lower edge of this empirical
band, which is the expected position once US shale entry is taken into
account: Huppmann's Figure~3 shows observed prices tracking the
calibrated oligopoly line below the pure-Cournot prediction, indicating
that OPEC behaves slightly less aggressively than pure Cournot would
imply. The qualitative finding that OPEC carries the largest Lerner in
every market structure is therefore robust across two independent
methodologies (ours, which is calibrated; theirs, which is
econometrically estimated) and is driven by the underlying cost
asymmetry rather than by any specific assumption about strategic
behaviour.

A Lerner index, however, only measures market power at the
simultaneous-move equilibrium. It is silent on whether the structural
cost advantage we have just documented can be amplified by strategic
timing. To answer that question, we now relax the simultaneity
assumption that Cournot imposes. The \citet[p.~??]{stackelberg1934} model
\citep[on whose original formulation see][]{baumgartner2001}
allows one player, which is the leader, to choose its output first; the
followers then observe the leader's quantity before choosing their own.
The leader can therefore optimise along the followers' best-response
functions rather than at the symmetric intersection point, which
generates a first-mover advantage \citep{dixit1980}\citep[p.~??]{tirole1988}. The size
of that advantage is an empirical question: it depends on both the cost
gap and the slope of the followers' reaction functions, and we now
quantify it for each candidate leader.

When OPEC acts as Stackelberg leader, it expands output to 93.33 mbd
(the largest single-player output in any structure), drives the price
down to \$46.67/bbl, and earns 2,133.33 in profit, a first-mover
advantage of 533.33 over the simultaneous Nash benchmark. When US shale
or Russia leads, the premium collapses to 75.00 and 208.33 respectively
(see Figure~6). The asymmetry has a clean economic reading: the leader's
benefit from expanding output rises with its cost advantage, because the
residual margin $P - c$ stays positive at quantities that drive
higher-cost rivals toward zero. \citet{beharritz2017} calibrate this
logic for the 2014 to 2016 episode and identify a market-share threshold
above which a low-cost swing producer prefers to flood the market rather
than to accommodate. Our Stackelberg-OPEC scenario, in which US profit
collapses from 225.00 to 1.67, is the quantitative counterpart of that
mechanism and reads naturally as the regime that Saudi Arabia
approximated during 2015 and again, more selectively, in March to April
2020. The combination of structural (cost) and strategic (timing)
advantage explains why no other producer has been able to displace OPEC
as the implicit residual setter of the OPEC+ quota envelope, despite
the substantial growth of US shale through 2024.

\subsection{Capacity constraints}
\label{subsec:capacity}

The preceding analysis assumes that each producer can supply any
quantity its best-response function calls for. In reality, each producer
faces a physical capacity ceiling: the upper bound on output that its
existing wells, pipelines, terminals, and processing infrastructure can
deliver in a given period. We calibrate $\mathrm{cap}_{\text{US}} = 30$,
$\mathrm{cap}_{\text{OPEC}} = 40$, and $\mathrm{cap}_{\text{RUS}} = 35$
mbd, levels consistent with current installed capacity plus a credible
six-month ramp; \appref{app:calibration} reports the empirical sources underlying
each value. A constraint is said to ``bind'' at a given equilibrium when
the player's optimal unconstrained output equals its cap; when this
happens, the cap restricts the strategic options of the binding player.
In a dynamic context, capacity is also a commitment device in the sense
of \citet{dixit1980}: unused capacity makes future flooding credible and
therefore alters rivals' incentives even when it is never deployed.

At the static Nash, OPEC's optimal unconstrained output of 40.00 mbd
exactly hits its capacity ceiling, so OPEC is the binding player; under
cartel quotas, no constraint binds, because the cooperative agreement
restricts output below every physical ceiling (see Figures~7 and~8).
Two implications follow. First, switching constraints on raises the
critical discount factor $\delta^*$ for every player, because binding
caps weaken the punishment phase by limiting the volume the cheated
party can flood with (see Figure~9). Second, expanding OPEC's capacity
ceiling has the opposite effect: sweeping
$\mathrm{cap}_{\text{OPEC}}$ from 25 to 60 mbd lowers $\delta^*$ from
0.850 to 0.529 (see Figure~10). Both statements are consistent because
they refer to different comparisons; we record this distinction
explicitly because it is easy to confuse.

The policy reading is direct. Saudi Arabia's persistent maintenance of
1.5 to 2.0 mbd of idle capacity, often criticized domestically as
wasteful, is rational because unused capacity is precisely what makes
the threat of retaliatory flooding credible and therefore lowers the
cooperation threshold for every member. \citet{till2016} in \appref{app:till}
provides the empirical counterpart, showing that prices become convex
in the supply deficit once OPEC spare capacity falls below a critical
threshold. The 2026 closure of the Strait of Hormuz on February~28
\citep{reuters2026}, which sent Brent above \$100/bbl and prompted a
record inventory drawdown \citep{iea2026omr}, is a contemporary
illustration of that convexity and of the value of spare capacity as a
stabilisation buffer.

\subsection{Welfare and deadweight loss}
\label{subsec:welfare}

We now evaluate the social impact of each market structure. To do so we
use the standard welfare decomposition of microeconomic theory. Consumer
surplus (CS) measures the area between the demand curve and the price
line, capturing the aggregate dollar value that consumers receive above
what they pay. Producer surplus (PS) measures the area between the
price line and the supply curve, capturing the aggregate margin that
producers earn above their marginal cost. Total welfare
$W = \mathrm{CS} + \mathrm{PS}$ is the social value generated by the
market. Deadweight loss (DWL) is the welfare lost relative to the
perfectly competitive benchmark, in which price equals marginal cost;
algebraically, $\mathrm{DWL} = W^{\text{competitive}} - W^{\text{actual}}$.
This decomposition follows \citet[p.~??]{mwg1995} and \citet[p.~??]{tirole1988}.

Under our calibration the competitive benchmark yields $W = 7{,}200$
with zero deadweight loss. At the static Nash, $\mathrm{DWL} = 800$ and
$W = 5{,}650$. Stackelberg-OPEC delivers $\mathrm{DWL} = 355.56$ and
$W = 6{,}627.78$, while the cartel inflicts $\mathrm{DWL} = 1{,}800$ and
$W = 4{,}837.50$, a welfare gap of $2{,}362.50$ units relative to
competition (see Figures~11 and~12). The deadweight-loss ranking, from
least to most distortionary, is competitive (0), Stackelberg-OPEC
(355.6), Stackelberg-Russia (501.4), Stackelberg-US (612.5), Nash
(800), and cartel (1,800). That Stackelberg leadership by the
lowest-cost player generates the highest welfare among realistic
oligopoly structures is counter-intuitive but economically clean: OPEC's
cost advantage means that production is allocated to the most efficient
producer, so the productive inefficiency of the symmetric Nash (in which
higher-cost firms produce too much) is reduced. The cartel is the worst
case because it combines allocative distortion (price above marginal
cost) with quota-induced productive distortion.

A particularly instructive extension is the interaction with a Pigouvian
per-barrel carbon tax $\tau$. The Pigouvian framework
\citep[p.~??]{pigou1920} prescribes a tax equal to the marginal external damage
caused by an activity, internalizing the externality at the producer
level. We sweep $\tau \in \{0, 10, 20, 30, 40, 50\}$. The collusion
premium, defined as the ratio of cartel to Nash joint profit minus one,
shrinks from 23.98 percent at $\tau = 0$ to 10.92 percent at $\tau = 50$
(see Figures~13 and~14). At high tax levels the cartel's output
restriction partially mimics the socially desired reduction in
fossil-fuel consumption, which we label the ``green-cartel'' effect:
collusion and climate policy point in the same direction. The result
should not be read as an endorsement of cartelization. A Pigouvian
carbon tax is strictly more efficient, because it achieves the same
quantity reduction while generating fiscal revenue and recycling part of
the loss into productive uses; cartels, in contrast, redistribute
roughly 1,400 units of consumer surplus to producers while destroying
1,000 units in deadweight loss.

\subsection{Bertrand price competition: robustness check}
\label{subsec:bertrand}

So far, we have assumed Cournot competition, in which producers choose
quantities. The natural alternative is Bertrand competition, in which
they choose prices. In its pure form, the \citet{bertrand1883} game with
homogeneous products and constant marginal cost produces the so-called
Bertrand paradox: even with only two producers, undercutting drives the
equilibrium price down to marginal cost. Pure-Bertrand is rarely a
realistic description of upstream oil markets because production is
capacity-constrained and crude grades are differentiated. We therefore
implement a differentiated-Bertrand specification in which each producer
faces a downward-sloping per-player demand and the substitution
parameter $\sigma \in [0,1]$ controls the degree of cross-grade
interchangeability: $\sigma = 0$ corresponds to independent monopolies
(no substitution), $\sigma = 1$ to perfect substitutes (the homogeneous
Bertrand case). This setting follows the \citet{maskintirole1988}
dynamic Bertrand framework and the standard differentiated-oligopoly
literature.

At the baseline $\sigma = 0.6$ (\appref{app:calibration} for reasoning behind this
calibration), which we read as the medium differentiation across WTI,
Brent, and Urals grades, the Bertrand-Nash average price is
\$35.90/bbl, far below the Cournot-Nash price of \$60.00/bbl (see
Figures~15 and~16). As \citet{bertrand1883} predicted, price competition
is more aggressive than quantity competition, and the equilibrium
converges toward marginal-cost pricing as $\sigma$ approaches one (see
Figure~17). Two qualifications must be flagged. The Bertrand and Cournot
specifications use different demand structures (per-player versus
aggregate), so absolute levels are not directly comparable; only the
qualitative ordering matters. We also document a corner-solution
artifact in the cooperative Bertrand optimisation at very low $\sigma$
(0.1 to 0.2), which we treat as a numerical limitation rather than an
economic finding.

Two qualitative results survive intact. The producer ranking
OPEC $>$ Russia $>$ US is preserved across competition modes, with
OPEC's profit rising in $\sigma$ because Bertrand convergence favours the
lowest-cost producer. The cooperation gap also persists: the cooperative
Bertrand average price exceeds the Bertrand-Nash by \$15.35, and joint
profit is roughly double (see Figure~18). This robustness, in the spirit
of \citet{maskintirole1988}, is reassuring: the central object of the
next section, the incentive to collude and the conditions under which
collusion can be sustained, does not depend on whether producers are
modelled as choosing quantities or prices. We retain Cournot for the
remainder of the analysis because it matches OPEC's institutional
reality, in which quotas, not posted prices, are the negotiated
instrument.
